\documentclass[12pt]{article}
\usepackage{amsmath,amssymb,amsfonts}
\usepackage[margin=1in]{geometry}
\begin{document}
\begin{center}
{\large\bfseries 17th Irish Mathematical Olympiad\\8 May 2004}
\end{center}
\bigskip\noindent\textbf{Paper 1}\bigskip
\begin{enumerate}
\item \begin{enumerate}
\item[(a)] For which positive integers $n$, does $2n$ divide the sum of the first $n$ positive integers?
\item[(b)] Determine, with proof, those positive integers $n$ (if any) which have the property that $2n + 1$ divides the sum of the first $n$ positive integers.
\end{enumerate}

\item Each of the players in a tennis tournament played one match against each of the others. If every player won at least one match, show that there is a group $A$, $B$, $C$ of three players for which $A$ beat $B$, $B$ beat $C$ and $C$ beat $A$.

\item $AB$ is a chord of length 6 of a circle centred at $O$ and of radius 5. Let $PQRS$ denote the square inscribed in the sector $OAB$ such that $P$ is on the radius $OA$, $S$ is on the radius $OB$ and $Q$ and $R$ are points on the arc of the circle between $A$ and $B$. Find the area of $PQRS$.

\item Prove that there are only two real numbers $x$ such that
\[
(x-1)(x-2)(x-3)(x-4)(x-5)(x-6) = 720.
\]

\item Let $a, b \geq 0$. Prove that
\[
\sqrt{2}\Bigl(\sqrt{a(a+b)^3} + b\sqrt{a^2+b^2}\Bigr) \leq 3(a^2 + b^2),
\]
with equality if and only if $a = b$.

\end{enumerate}
\newpage
\begin{center}
{\large\bfseries 17th Irish Mathematical Olympiad\\8 May 2004}
\end{center}
\bigskip\noindent\textbf{Paper 2}\bigskip
\begin{enumerate}
\setcounter{enumi}{5}
\item Determine all pairs of prime numbers $(p, q)$, with $2 \leq p, q < 100$, such that $p + 6$, $p + 10$, $q + 4$, $q + 10$ and $p + q + 1$ are all prime numbers.

\item $A$ and $B$ are distinct points on a circle $T$. $C$ is a point distinct from $B$ such that $|AB| = |AC|$, and such that $BC$ is tangent to $T$ at $B$. Suppose that the bisector of $\angle ABC$ meets $AC$ at a point $D$ inside $T$. Show that $\angle ABC > 72^\circ$.

\item Suppose $n$ is an integer $\geq 2$. Determine the first digit after the decimal point in the decimal expansion of the number
\[
\sqrt[3]{n^3 + 2n^2 + n}.
\]

\item Define the function $m$ of the three real variables $x, y, z$ by
\[
m(x, y, z) = \max(x^2, y^2, z^2), \quad x, y, z \in \mathbb{R}.
\]
Determine, with proof, the minimum value of $m$ if $x, y, z$ vary in $\mathbb{R}$ subject to the following restrictions:
\[
x + y + z = 0, \quad x^2 + y^2 + z^2 = 1.
\]

\item Suppose $p$, $q$ are distinct primes and $S$ is a subset of $\{1, 2, \ldots, p-1\}$. Let $N(S)$ denote the number of solutions of the equation
\[
\sum_{i=1}^{q} x_i \equiv 0 \pmod{p},
\]
where $x_i \in S$, $i = 1, 2, \ldots, q$. Prove that $N(S)$ is a multiple of $q$.

\end{enumerate}
\end{document}
